\begin{theorem}[Bernoulli’s inequality]
Let $x\in\mathbb R$ satisfy $x\ge -1$ and let $n\in\mathbb Z_{\ge0}$. Then
\[
(1+x)^{\,n}\ge 1+n\,x .
\]
\end{theorem}

\begin{proof}
\textbf{1. Degenerate and trivial cases.}\\
\emph{Case $n=0$.} $(1+x)^{0}=1=1+0\cdot x$.\\
\emph{Case $x=0$.} $(1+0)^{n}=1=1+n\cdot0$.\\
\emph{Case $x=-1$.} $(1+x)^{n}=0$, while $1+n x = 1-n$.  
If $n=1$ we have $0=0$; if $n\ge2$, then $1-n\le -1$, so $0\ge 1-n$.  
Thus the inequality holds for the endpoint $x=-1$.

Henceforth assume $x>-1$ (so $1+x>0$) and $n\ge1$.

\medskip\noindent\textbf{2. Basis of induction.}\\
For $n=1$,
\[
(1+x)^{1}=1+x=1+1\cdot x ,
\]
so the statement is true (with equality).

\medskip\noindent\textbf{3. Inductive hypothesis.}\\
Assume that for a fixed $k\ge1$ we have
\begin{equation}
(1+x)^{k}\ge 1+kx \qquad\text{for all }x\ge-1 .
\tag{IH}
\end{equation}

\medskip\noindent\textbf{4. Inductive step.}\\
Multiplying \eqref{IH} by the non‑negative factor $1+x$ yields
\begin{equation}
(1+x)^{k+1}=(1+x)^{k}(1+x)\ge (1+kx)(1+x) .
\tag{1}
\end{equation}
Expanding the right‑hand side,
\begin{align}
(1+kx)(1+x) &= 1 + (k+1)x + kx^{2} .
\tag{2}
\end{align}
Substituting \eqref{2} into \eqref{1} gives
\begin{equation}
(1+x)^{k+1}\ge 1+(k+1)x+kx^{2} .
\tag{3}
\end{equation}
Since $x^{2}\ge0$ and $k\ge1$, the term $k x^{2}$ is non‑negative; discarding it preserves the inequality:
\[
(1+x)^{k+1}\ge 1+(k+1)x .
\]
Thus the statement holds for $n=k+1$.

\medskip\noindent\textbf{5. Completion of the induction.}\\
The inequality is true for $n=1$ (basis) and the inductive step shows that truth for $n=k$ implies truth for $n=k+1$. By the principle of mathematical induction it holds for all integers $n\ge1$. Together with the trivial case $n=0$ proved in step 1, the theorem is valid for every $n\ge0$.

\medskip\noindent\textbf{6. Equality cases.}\\
From \eqref{3} we have
\[
(1+x)^{k+1}=1+(k+1)x \iff kx^{2}=0 .
\]
Because $k\ge1$, this forces $x=0$.  
Consequently:
\begin{itemize}
\item for $n=0$ equality holds for every admissible $x$;
\item for $n=1$ equality holds for every admissible $x$;
\item for $n\ge2$ equality holds \emph{iff} $x=0$.
\end{itemize}
In compact form,
\[
(1+x)^{n}=1+n x \quad\Longleftrightarrow\quad
\begin{cases}
x=0,\; \text{any }n\ge0,\\[2pt]
\text{or } n\in\{0,1\},\; \text{any }x\ge-1 .
\end{cases}
\]

\qed
\end{proof}